\documentclass[conference]{IEEEtran}
\IEEEoverridecommandlockouts

%% ---- Core packages ----
\usepackage{cite}
\usepackage{amsmath,amssymb,amsfonts}
\usepackage{algorithmic}
\usepackage{graphicx}
\usepackage{textcomp}
\usepackage{xcolor}
\usepackage{booktabs}
\usepackage{multirow}
\usepackage{array}
\usepackage{url}
\usepackage{hyperref}
\usepackage{listings}
\usepackage{balance}

%% ---- hyperref setup ----
\hypersetup{
  colorlinks=true,
  linkcolor=blue,
  citecolor=blue,
  urlcolor=blue,
  pdftitle={SVAC-Concurrency: A Step-Verification Benchmark for LLM Reasoning on OS Concurrency Algorithms},
  pdfauthor={[Author Name]},
  pdfkeywords={Large Language Models, Algorithmic Reasoning, Operating Systems, Concurrency Control,
               Deadlock Detection, Banker's Algorithm, Step-level Verification, Benchmark}
}

%% ---- Custom column types for tables ----
\newcolumntype{L}[1]{>{\raggedright\arraybackslash}p{#1}}
\newcolumntype{C}[1]{>{\centering\arraybackslash}p{#1}}
\newcolumntype{R}[1]{>{\raggedleft\arraybackslash}p{#1}}

%% ---- Convenience macros ----
\newcommand{\SA}{\text{SA}}
\newcommand{\FEP}{\text{FEP}}
\newcommand{\ECR}{\text{ECR}}
\newcommand{\VA}{\text{VA}}
\newcommand{\SAvalid}{\text{SA}_{\text{valid}}}

% ============================================================
\begin{document}
% ============================================================

\title{SVAC-Concurrency: A Step-Verification Benchmark for Large Language Model
       Reasoning on Operating System Concurrency Algorithms}

%% Replace the author block with real names before submission.
\author{%
  \IEEEauthorblockN{[Author Name]\textsuperscript{1},
                    [Advisor / Co-Author Name]\textsuperscript{1}}
  \IEEEauthorblockA{\textsuperscript{1}Department of Computer Science and Engineering,\\
    [Institution / University Name], [City, State, Country]\\
    Email: \{author.name, advisor.name\}@university.edu}
}

\maketitle

% ============================================================
\begin{abstract}
% ============================================================
While Large Language Models (LLMs) demonstrate remarkable competence in natural language
understanding and high-level source code synthesis, their ability to faithfully track
deterministic, multi-step algorithmic state transitions remains fundamentally unverified.
Concurrency control and resource allocation algorithms in Operating Systems (OS) represent
a rigorous, safety-critical testbed for mechanistic reasoning, demanding strict invariant
preservation, vector/matrix state updates, and topological cycle detection.

In this paper, we present \textbf{SVAC-Concurrency}, an open-source, reproducible
evaluation benchmark comprising \textbf{125 parameterized problem instances} (generating
375 evaluation traces per model across 3 prompting regimes) spanning five canonical OS
algorithms: Wait-For Graph (WFG) deadlock detection, Banker's safety algorithm, Buddy
System memory allocation, counting Semaphores, and Mutex lock simulations.

Rather than relying solely on coarse outcome correctness (Verdict Accuracy),
SVAC-Concurrency introduces a deterministic ground-truth verification framework to
evaluate execution traces across four diagnostic dimensions: \textbf{Step Accuracy (SA)},
\textbf{First Error Position (FEP)}, \textbf{Error Cascade Rate (ECR)}, and
\textbf{Equivalence-Aware Invariant Validity}. We conduct an extensive cross-architectural
investigation evaluating Google Gemini-3.5-Flash-Lite, OpenAI GPT-OSS-20B, NVIDIA
Nemotron-3-Super-120B, and Liquid AI LFM-2.6B across Zero-Shot,
Few-Shot, and Chain-of-Thought (CoT) prompting strategies.

Our empirical findings reveal:
\begin{enumerate}
\item \textbf{Procedural Fidelity vs.\ Outcome Illusion:} Across balanced 50/50 decision
  distributions (majority baseline $= 50.0\%$), Gemini-3.5-Flash-Lite achieves 72.8\%
  Verdict Accuracy on parsed responses but only 56.1\% Step Accuracy---collapsing to
  \textbf{36.98\% SA} on non-saturating state-tracking tasks (WFG, Banker's, Buddy)---
  demonstrating widespread ``right verdict for flawed reasons'' behavior.
\item \textbf{Prompt-Priming vs.\ Invariant Divergence in Graph Reasoning:} Under strict
  positional matching, WFG cycle detection exhibits an apparent collapse from 85.4\% SA
  in Few-Shot to 7.5\% in CoT (and 4.4\% in Zero-Shot). Crucially,
  equivalence-aware invariant verification reveals that CoT maintains
  \textbf{93.2\% valid step accuracy} and \textbf{93.3\% verdict accuracy} (comparable
  to Few-Shot's 100.0\% valid SA and 100.0\% VA). The collapse in strict positional SA
  is driven by prompt-primed canonical tie-breaking rather than algorithmic incapacity,
  demonstrating that rigid sequential benchmarks risk measuring format compliance rather
  than reasoning.
\item \textbf{Failure Horizon Dichotomy:} WFG errors under strict matching occur at Step
  1 in 100\% of Zero-Shot and CoT failure cases due to branch tie-breaking divergence,
  whereas Buddy allocation exhibits true intermediate state corruption (mean
  $\FEP = 6.6$--$7.6$ steps).
\item \textbf{Equivalence vs.\ Strict Positional Bias:} Invariant-based scoring reveals
  that strict positional matching severely underestimates model competence: valid step
  accuracy increases to 74.3\% in Banker's ($\Delta = {+}32.8\%$) and 95.6\% in WFG
  ($\Delta = {+}64.3\%$).
\item \textbf{Format Brittleness in Open Architectures:} Severe schema compliance
  degradation in open-weight models (19.7\% parse success for Nemotron-120B) demonstrates
  that raw parameter scale does not guarantee structured state-machine adherence.
\end{enumerate}

The complete benchmark suite, deterministic reference solvers, prompt templates, and
evaluation pipelines are publicly available at:
\url{https://github.com/anonymous-author/SVAC-Concurrency}.
\end{abstract}

\begin{IEEEkeywords}
Large Language Models, Algorithmic Reasoning, Operating Systems, Concurrency Control,
Deadlock Detection, Banker's Algorithm, Step-level Verification, Benchmark.
\end{IEEEkeywords}

% ============================================================
\section{Introduction}
\label{sec:intro}
% ============================================================

Large Language Models (LLMs) are increasingly integrated into modern software engineering
pipelines, automated code review tools, and operating systems pedagogy~\cite{brown2020language,bubeck2023sparks}.
Despite remarkable empirical gains on functional coding benchmarks such as
HumanEval~\cite{chen2021codex}, MBPP~\cite{austin2021program}, and
APPS~\cite{hendrycks2021apps}, conventional benchmarks evaluate programs almost
exclusively via black-box input-output test cases. They fail to assess whether an
auto-regressive model maintains a coherent, mechanically faithful simulation of
intermediate program states during execution~\cite{cobbe2021training,hendrycks2021math,lightman2024verify}.

In operating systems and systems programming, procedural correctness is
non-negotiable. Concurrency control algorithms---such as deadlock detection, resource
allocation safety, and dynamic physical memory partitioning---are governed by rigorous
theoretical properties:
\begin{enumerate}
\item \textbf{Strict Discrete State Tracking:} Every transition must update resource
  availability vectors, process allocation matrices, or block allocation trees without
  state drift or phantom resource generation.
\item \textbf{Topological and Relational Invariants:} Identifying deadlocks requires
  systematic traversal of directed Wait-For Graphs to detect closed cycles without
  hallucinating non-existent dependency edges.
\item \textbf{Catastrophic Error Cascades:} A single state corruption at step $t$
  invalidates all subsequent state assertions at $t{+}k$, converting downstream
  execution into an unfaithful hallucination.
\end{enumerate}

When evaluating LLMs on such procedural tasks, standard outcome-based metrics (e.g.,
binary classification of whether a system is deadlocked or safe) introduce massive
evaluation noise. On balanced datasets, an LLM can easily guess the correct binary
verdict through statistical heuristics or shallow token correlations while executing an
entirely invalid sequence of intermediate states.

To resolve this evaluation gap, we introduce \textbf{SVAC-Concurrency}
(\emph{Step-Verification of Algorithmic Concurrency}), a comprehensive,
parameter-controlled benchmark designed to evaluate intermediate LLM reasoning traces
across five foundational OS concurrency algorithms. Each instance is paired with an
exact ground-truth execution trace computed by deterministic reference solvers,
unit-tested against textbook worked examples~\cite{silberschatz2018os} and, for WFG
and Banker's algorithms (60 of 125 instances), validated via independent brute-force
differential testing.

\subsection*{Research Questions}
\begin{itemize}
\item \textbf{RQ1} (\emph{Procedural Fidelity vs.\ Outcome Correctness}): Do LLMs
  produce faithful intermediate state transitions, or do high verdict accuracies mask
  severe intermediate state-tracking failures on balanced decision distributions?
\item \textbf{RQ2} (\emph{Complexity Scaling \& Failure Horizons}): How do increasing
  state-space bounds affect step accuracy, and at what execution horizon ($\FEP$) do
  cascading errors initiate across different algorithmic structures?
\item \textbf{RQ3} (\emph{Prompting Dynamics \& Serialization Artifacts}): Does
  step-by-step verbalization (CoT) assist or destabilize discrete graph-traversal and
  matrix-tracking algorithms, and to what extent does few-shot prompting act as format
  serialization priming?
\item \textbf{RQ4} (\emph{Equivalence \& Architectural Robustness}): To what extent does
  strict positional scoring underestimate valid alternative execution paths, and how do
  open-weight models adhere to strict intermediate state schemas compared to
  instruction-tuned frontier APIs?
\end{itemize}

\subsection*{Primary Contributions}
\begin{itemize}
\item \textbf{Comprehensive Benchmark Suite:} 125 unique instances across 5 core OS
  algorithms parameterized across 5 distinct complexity tiers ($N{=}375$ evaluations
  per model across 3 prompt strategies).
\item \textbf{Multi-Dimensional Metric Taxonomy:} Formalization of $\SA$, $\FEP$,
  $\ECR$, and Equivalence-Aware Invariant Validity to capture reasoning dynamics beyond
  binary verdicts.
\item \textbf{Cross-Architectural Empirical Investigation:} Systematic comparison of four
  distinct model families: Gemini-3.5-Flash-Lite, GPT-OSS-20B, Nemotron-3-Super-120B,
  and Liquid AI LFM-2.6B.
\item \textbf{Deconstruction of the CoT Inversion Effect:} Discovery and formal
  characterization showing that the apparent collapse of CoT in graph cycle detection
  (85.4\% Few-Shot vs.\ 7.5\% CoT strict SA) is primarily driven by serialization
  tie-breaking divergence, while invariant-aware step validity remains robust at 93.2\%
  with 93.3\% verdict accuracy.
\item \textbf{Equivalence-Aware Invariant Evaluation:} Demonstration that rigid
  positional matching under-reports valid algorithmic execution by 32.8\% in Banker's
  and 64.3\% in WFG.
\end{itemize}

% ============================================================
\section{Related Work}
\label{sec:related}
% ============================================================

\subsection{Benchmarking Code Generation and Execution Reasoning}

Early LLM code evaluation benchmarks such as HumanEval~\cite{chen2021codex},
MBPP~\cite{austin2021program}, and APPS~\cite{hendrycks2021apps} established
input-output testing as the standard paradigm for evaluating program synthesis. Recent
efforts have shifted toward evaluating \emph{execution reasoning}. CRUXEval~\cite{gu2024cruxeval}
evaluated code execution reasoning by testing input/output prediction across short Python
snippets. REval~\cite{chen2025reval} and CodeMind~\cite{liu2024codemind} introduced
state-level execution tracing for general imperative code. While these benchmarks test
arbitrary code snippets, they do not systematically probe canonical discrete algorithms
where formal invariants, resource constraints, and topological cycles dictate
correctness. SVAC-Concurrency provides precisely this specialized, parameter-controlled
algorithmic framework.

\subsection{Mechanistic State Tracking in Transformers}

A growing body of literature investigates whether transformer architectures maintain
coherent internal world models during sequential tasks. Toshniwal et al.~\cite{toshniwal2022chess}
demonstrated that auto-regressive language models track board state representations in
chess games despite receiving only move sequences. Strobl et al.~\cite{strobl2024transformers}
surveyed the theoretical boundaries of what formal languages and state machines
transformer architectures can express. Wang et al.~\cite{wang2023graph} investigated
graph traversal algorithms in natural language. In the systems domain, the precursor SVAC
technical report~\cite{anon2024svac} evaluated step-level reasoning in virtual memory
page replacement policies (FIFO, LRU, Optimal, Clock). Additionally, the CONCUR
benchmark by Huang et al.~\cite{huang2026concur} evaluated LLMs on \emph{synthesizing}
concurrent synchronization primitives. Our work explicitly contrasts with CONCUR: rather
than evaluating concurrent code synthesis, we evaluate \emph{execution-trace
verification}, isolating whether models accurately track concurrent state transitions and
deadlock invariants.

\subsection{Prompting Dynamics and CoT Unfaithfulness}

Chain-of-Thought (CoT) prompting~\cite{wei2022chain,chu2024survey} is widely documented to boost
arithmetic and multi-hop reasoning. However, recent empirical studies reveal significant
limits and pathology in verbalized reasoning. Lanham et al.~\cite{lanham2023faithfulness}
and Turpin et al.~\cite{turpin2023unfaithful} demonstrated that CoT explanations are
frequently unfaithful, serving as post-hoc rationalizations rather than causal
computation drivers. Shi et al.~\cite{shi2023distracted} demonstrated that LLMs are
easily distracted by irrelevant reasoning context, while Liu et al.~\cite{liu2024mindstep}
observed degradation in multi-step planning tasks when natural language narration exceeds
formal state-space capacity. In graph-theoretic tasks, Agrawal et al.~\cite{agrawal2024graph}
investigated the structural limitations of language models on graph reasoning. Our work
grounds these phenomena in operating systems concurrency, providing quantitative evidence
of how serialization tie-breaks and verbalized narration interact with discrete state tracking.

% ============================================================
\section{Benchmark Architecture and Methodology}
\label{sec:benchmark}
% ============================================================

\begin{table*}[t]
\centering
\caption{SVAC-Concurrency Benchmark Suite Overview}
\label{tab:suite}
\begin{tabular}{L{2.8cm}L{3.2cm}C{1.4cm}L{3.4cm}L{3.0cm}}
\toprule
\textbf{Algorithm Family} & \textbf{State Representation} & \textbf{Instances} &
\textbf{Complexity Bounds} & \textbf{Class Balance} \\
\midrule
Wait-For Graph (WFG) & Directed Adjacency \& Recursion Stack & 30 &
  $|V| \in [2,6]$, $|E| \in [1,8]$ & 50\% Deadlock (15) / 50\% Acyclic (15) \\
\addlinespace
Banker's Algorithm & Available Vector, Allocation/Need Matrix & 30 &
  $P \in [2,6]$, $R \in [2,4]$ & 50\% Safe (15) / 50\% Unsafe (15) \\
\addlinespace
Buddy System Memory & Free-list Binary Tree \& Blocks & 25 &
  512\,KB--4096\,KB, Requests $\in [4,10]$ & Deterministic Partition Trace \\
\addlinespace
Counting Semaphore & \texttt{S.value} \& FIFO Blocked Queue & 20 &
  $S \in [1,3]$, Ops $\in [4,10]$ & Deterministic State Sequence \\
\addlinespace
Mutex Lock Synchronization & Owner ID \& Lock State Flag & 20 &
  Threads $\in [2,5]$, Ops $\in [4,8]$ & Deterministic Transition Trace \\
\midrule
\textbf{Total} & & \textbf{125} & 5 Complexity Tiers & 50.0\% Majority Baseline \\
\bottomrule
\end{tabular}
\end{table*}

\subsection{Problem Formulations \& Ground-Truth Verification}

All problem instances are deterministically synthesized across five calibrated complexity
tiers (Levels 1--5). Ground-truth execution traces are produced by reference Python
solvers (\texttt{solvers/}) unit-tested against textbook worked examples from Silberschatz
et al.~\cite{silberschatz2018os}. Additionally, WFG and Banker's Algorithm solvers
(60 of 125 instances) were independently cross-verified via brute-force differential
testing against alternative reference implementations (3-color DFS oracle for WFG;
exhaustive permutation search for Banker's).

\begin{enumerate}
\item \textbf{Wait-For Graph (WFG) Deadlock Detection:} Given a set of processes
  $P = \{P_0, \dots, P_{n-1}\}$ and directed wait edges $E \subseteq P \times P$, where
  $(P_i, P_j) \in E$ indicates $P_i$ is blocked waiting for a resource held by $P_j$.
  The model must execute DFS, tracking the active recursion stack and visited set at each
  step, identifying back-edges, reporting the cycle subgraph if present, and emitting a
  final verdict $\text{deadlock} \in \{\text{true}, \text{false}\}$. The dataset is
  strictly balanced: 15 instances contain cycles and 15 are directed acyclic graphs
  (DAGs), balanced evenly (3 positive, 3 negative) across all 5 complexity tiers.

\item \textbf{Banker's Algorithm for Deadlock Avoidance:} Given an Available resource
  vector $A \in \mathbb{N}^m$, Allocation matrix $\mathbf{Alloc} \in \mathbb{N}^{n \times m}$,
  and Max demand matrix $\mathbf{Max} \in \mathbb{N}^{n \times m}$. The model must
  compute $\mathbf{Need} = \mathbf{Max} - \mathbf{Alloc}$, iteratively identify an
  executable candidate process $P_i$ such that $\mathbf{Need}_i \le \text{Work}$,
  simulate resource reclamation $\text{Work} \leftarrow \text{Work} + \mathbf{Alloc}_i$,
  update the $\text{Finish}$ boolean array, and determine whether a complete safe
  sequence exists ($\text{safe} \in \{\text{true}, \text{false}\}$). Strictly balanced:
  15 safe and 15 unsafe instances (3 per complexity tier).

\item \textbf{Buddy System Memory Allocation:} Given total memory size $M = 2^K$ and a
  sequential stream of allocation and deallocation operations. The model must track
  recursive block splits, power-of-two alignments, and buddy coalescing on free
  operations, returning the exact block partition state at each step.

\item \textbf{Counting Semaphore Simulation:} Given initial counter $S \ge 0$ and an
  interleaved execution stream of concurrent $\text{wait}(S)$ and $\text{signal}(S)$
  calls from multiple threads. The model must track the integer value of $S$ and the FIFO
  blocked thread queue at every transition.

\item \textbf{Mutex Lock Simulation:} Evaluates mutual exclusion lock acquisition,
  priority ownership, contention blocking, and release transitions across concurrent
  threads.
\end{enumerate}

\subsection{Mathematical Evaluation Framework}

Let an evaluation instance $I$ have ground-truth trace
$\mathbf{Y}^* = (y_1^*, y_2^*, \dots, y_{T^*}^*)$ of length $T^*$, and model-generated
trace $\hat{\mathbf{Y}} = (\hat{y}_1, \hat{y}_2, \dots, \hat{y}_m)$ of length $m$.

\subsubsection{Formal Step Equality Specification}

A step $\hat{y}_t$ matches ground truth $y_t^*$ if and only if all task-specific state
attributes match under canonical normalization:

\textbf{WFG Deadlock:}
\begin{align}
\hat{y}_t = y_t^* \iff
  &\bigl(\hat{y}_t.\text{action} = y_t^*.\text{action}\bigr) \notag\\
  &\land \bigl(\hat{y}_t.\text{node} = y_t^*.\text{node}\bigr) \notag\\
  &\land \bigl(\mathrm{sort}(\hat{y}_t.\text{visited}) = \mathrm{sort}(y_t^*.\text{visited})\bigr)\notag\\
  &\land \bigl(\hat{y}_t.\text{stack} = y_t^*.\text{stack}\bigr)
\end{align}
Visited nodes are compared as canonical sorted sets (order-insensitive), while the
recursion stack is compared as an ordered sequence (order-sensitive).

\textbf{Banker's Algorithm:}
\begin{align}
\hat{y}_t = y_t^* \iff
  &\bigl(\hat{y}_t.\text{candidate} = y_t^*.\text{candidate}\bigr) \notag\\
  &\land \bigl(\hat{y}_t.\text{work\_before} = y_t^*.\text{work\_before}\bigr)\notag\\
  &\land \bigl(\hat{y}_t.\text{work\_after} = y_t^*.\text{work\_after}\bigr)\notag\\
  &\land \bigl(\hat{y}_t.\text{finish} = y_t^*.\text{finish}\bigr)
\end{align}

\textbf{Buddy System / Semaphore / Mutex:} Exact equality across all state fields
(allocated block address-size tuples, free-list, counter/owner state, blocked thread
list order).

\subsubsection{Diagnostic Trace Metrics}

\begin{enumerate}
\item \textbf{Step Accuracy ($\SA$):}
  \begin{equation}
    \SA = \frac{1}{T^*} \sum_{t=1}^{\min(T^*,m)} \mathbb{I}(\hat{y}_t = y_t^*)
  \end{equation}

\item \textbf{First Error Position ($\FEP$):}
  \begin{equation}
    \FEP = \min\!\left(\bigl\{t \in \{1,\dots,T^*\} \mid \hat{y}_t \ne y_t^*\bigr\}
             \cup \{\infty\}\right)
  \end{equation}

\item \textbf{Error Cascade Rate ($\ECR$):}
  \begin{equation}
    \ECR = \begin{cases}
      \dfrac{1}{T^* - \FEP} \displaystyle\sum_{t=\FEP+1}^{T^*}
        \mathbb{I}(\hat{y}_t \ne y_t^*), & \text{if } \FEP < T^* \\[6pt]
      0, & \text{if } \FEP \ge T^* \text{ or } \FEP = \infty
    \end{cases}
  \end{equation}

\item \textbf{Verdict Accuracy ($\VA$):} To prevent conflating parse failures with
  reasoning errors, we report both $\VA_{\text{parsed}}$ (conditionally over successfully
  formatted outputs) and $\VA_{\text{all}}$ (treating unparsed outputs as incorrect):
  \begin{align}
    \VA_{\text{parsed}} &= \frac{1}{N_{\text{parsed}}}
      \sum_{i=1}^{N_{\text{parsed}}} \mathbb{I}(\widehat{\text{verdict}}_i = \text{verdict}_i^*)\\
    \VA_{\text{all}} &= \frac{1}{N_{\text{total}}}
      \sum_{i=1}^{N_{\text{total}}} \mathbb{I}\!\left(\text{parsed}_i \land
      \widehat{\text{verdict}}_i = \text{verdict}_i^*\right)
  \end{align}
\end{enumerate}

\subsubsection{Equivalence-Aware Invariant Scoring}

Strict positional matching penalizes models when multiple valid execution paths exist
(e.g., Banker's instances where multiple processes satisfy $\mathbf{Need}_i \le \text{Work}$,
or graph traversals with alternative neighbor ordering). We therefore formulate
\textbf{Equivalence-Aware Invariant Validity ($\SAvalid$)}:
\begin{itemize}
\item At each step $t$, verify that $\hat{y}_t.\text{candidate}$ is unfinished and
  satisfies $\mathbf{Need}_{\text{candidate}} \le \text{Work}_{t-1}$.
\item Verify that $\text{Work}_t = \text{Work}_{t-1} + \mathbf{Alloc}_{\text{candidate}}$
  and only the candidate process finishes.
\item If safe, verify that $\hat{y}.\text{safe\_sequence}$ is an executable permutation
  under initial resources.
\end{itemize}

% ============================================================
\section{Experimental Setup}
\label{sec:setup}
% ============================================================

All evaluations were executed under strict deterministic parameters
($\text{temperature} = 0.0$, $\text{top\_p} = 1.0$) across three prompt configurations:
\textbf{Zero-Shot (ZS)} (direct JSON schema specification), \textbf{Few-Shot (FS)}
(in-context exemplars of step-by-step execution), and \textbf{Chain-of-Thought (CoT)}
(verbalized rationale preceding structured JSON emission).

We evaluated four diverse foundation models:
\begin{enumerate}
\item \textbf{Google Gemini-3.5-Flash-Lite:} High-throughput dense instruction-tuned
  reasoning model evaluated via Google AI Studio API endpoints across all 125 instances
  ($N{=}375$).
\item \textbf{OpenAI GPT-OSS-20B:} Open-weight model evaluated via Groq high-speed LPU
  inference acceleration, scoped to WFG deadlock detection ($N{=}81$ scored attempts) due
  to upstream API rate-limits.
\item \textbf{NVIDIA Nemotron-3-Super-120B:} 120B-parameter open-weight foundation model
  evaluated via OpenRouter endpoints across all 125 instances ($N{=}375$ attempts).
\item \textbf{Liquid AI LFM-2.6B:} Non-transformer dynamical state-space architecture
  evaluated via OpenRouter endpoints, scoped to WFG ($N{=}39$ scored attempts).
\end{enumerate}

% ============================================================
\section{Results and Empirical Analysis}
\label{sec:results}
% ============================================================

\begin{table*}[t]
\centering
\caption{Multi-Provider Cross-Architecture Benchmark Results}
\label{tab:cross_model}
\begin{tabular}{L{2.8cm}C{2.0cm}C{1.4cm}C{1.4cm}C{1.4cm}C{1.4cm}C{1.7cm}C{1.7cm}}
\toprule
\textbf{Model Architecture} & \textbf{Organization} & \textbf{Att.} & \textbf{Parse}
  & \textbf{SA} & \textbf{ECR} & \textbf{VA\textsubscript{parsed}} & \textbf{VA\textsubscript{all}} \\
\midrule
Gemini-3.5-Flash-Lite & Google DeepMind & 375 & 99.2\% & 56.1\% & 36.9\% & 72.8\% & 72.3\% \\
GPT-OSS-20B           & OpenAI          & 81\rlap{$^\dagger$} & 58.0\% & 53.1\%$^*$ & 45.3\%$^*$ & 100.0\%$^*$ & 58.0\% \\
Nemotron-3-Super-120B & NVIDIA          & 375 & 19.7\% & 48.5\%$^*$ & 41.3\%$^*$ & 95.9\%$^*$ & 18.9\% \\
Liquid LFM-2.6B       & Liquid AI       & 39\rlap{$^\dagger$} & 69.2\% & 14.4\%$^*$ & 86.1\%$^*$ & 96.3\%$^*$ & 66.7\% \\
\midrule
Majority Baseline & --- & 375 & 100.0\% & 0.0\% & 100.0\% & 50.0\% & 50.0\% \\
\bottomrule
\multicolumn{8}{l}{\footnotesize $^*$Computed conditionally over successfully parsed JSON outputs.}\\
\multicolumn{8}{l}{\footnotesize $^\dagger$Scoped to WFG subset due to upstream provider rate-limiting and connection resets.}\\
\multicolumn{8}{l}{\footnotesize Parse Rate $= N_{\text{parsed}} / N_{\text{attempted}}$. Att.~= attempts.}
\end{tabular}
\end{table*}

\begin{table}[t]
\centering
\caption{Per-Task Step Accuracy (\%) Across Prompt Strategies\\
         (Gemini-3.5-Flash-Lite, $N = 20$--$30$ instances per cell)}
\label{tab:task_sa}
\begin{tabular}{L{2.4cm}C{1.4cm}C{1.4cm}C{1.4cm}}
\toprule
\textbf{Task Domain} & \textbf{ZS} & \textbf{FS} & \textbf{CoT} \\
\midrule
WFG Deadlock         & 4.4\%  & 85.4\% & 7.5\%  \\
Banker's Safety      & 32.9\% & 41.9\% & 49.7\% \\
Buddy System         & 30.5\% & 31.5\% & 49.1\% \\
Counting Semaphore   & 91.2\% & 96.2\% & 95.0\% \\
Mutex Lock           & 100.0\%& 100.0\%& 100.0\%\\
\midrule
Non-Saturated Mean\rlap{$^\ddagger$} & 22.6\% & 52.9\% & 35.4\% \\
Instance-Weighted Mean\rlap{$^\natural$} & 45.6\% & 68.2\% & 54.7\% \\
\bottomrule
\multicolumn{4}{l}{\footnotesize $^\ddagger$WFG + Banker's + Buddy only. Overall non-saturated SA: 36.98\%.}\\
\multicolumn{4}{l}{\footnotesize $^\natural$Weighted by $N$ per task (30,30,25,20,20); strategy mean = 56.2\%, reconciling to suite SA of 56.1\%.}
\end{tabular}
\end{table}

\subsection{The ``Right Verdict for Wrong Reasons'' Gap}

A primary finding from our benchmark is the stark divergence between Verdict Accuracy and
Step Accuracy. As shown in Table~\ref{tab:cross_model}, Gemini-3.5-Flash-Lite achieves an
overall \textbf{Verdict Accuracy of 72.8\%} on parsed responses, yet exhibits a mean
\textbf{Step Accuracy of only 56.1\%}.

Crucially, the 56.1\% suite mean is heavily inflated by trivial state saturation in
Mutex (100.0\% SA) and Semaphore tasks (94.2\% SA). When isolating the \textbf{non-saturating
algorithmic tasks (WFG, Banker's, and Buddy System)}, mean Step Accuracy drops to
\textbf{36.98\%}. Because the decision tasks are strictly 50/50 balanced (majority
baseline $= 50.0\%$), the gap between 72.8\% VA and 37.0\% non-saturated SA demonstrates
that outcome accuracy vastly overstates procedural competence.

\subsection{Deconstructing the CoT Inversion Effect: Priming vs.\ Invariants}

Table~\ref{tab:task_sa} appears to show an extreme prompting anomaly. In linear-procedural
tasks, CoT prompting provides substantial gains ($+16.8\%$ Banker's, $+18.6\%$ Buddy
over Zero-Shot). In WFG cycle detection, strict positional matching reports an apparent
collapse: Few-Shot attains \textbf{85.4\% SA} whereas CoT plunges to \textbf{7.5\% SA}.

\begin{table}[t]
\centering
\caption{WFG Empirical Error Taxonomy \& Accuracy Breakdown (Gemini)}
\label{tab:wfg_taxonomy}
\begin{tabular}{L{3.2cm}C{1.5cm}C{1.5cm}C{1.3cm}}
\toprule
\textbf{Metric / Error Category} & \textbf{ZS} ($N{=}30$) & \textbf{FS} ($N{=}30$) & \textbf{CoT} ($N{=}30$) \\
\midrule
Visited Set Corruption    & 30 (100\%)  &  7 (23.3\%) & 30 (100\%) \\
Recursion Stack Corr.     & 14 (46.7\%) & 12 (40.0\%) &  8 (26.7\%) \\
Wrong Traversal Node      & 20 (66.7\%) & 12 (40.0\%) & 18 (60.0\%) \\
Step Count Mismatch       & 16 (53.3\%) & 13 (43.3\%) & 12 (40.0\%) \\
False Cycle Hallucination &  3 (10.0\%) &  0 (0.0\%)  &  2 (6.7\%)  \\
Missed Cycle (False Neg.) &  0 (0.0\%)  &  0 (0.0\%)  &  0 (0.0\%)  \\
\midrule
Verdict Acc.\ (VA)        & 27/30 (90.0\%) & 28/28 (100\%)\rlap{$^\S$} & 28/30 (93.3\%) \\
Strict Positional SA      &  4.4\%  & 85.4\% &  7.5\%  \\
Invariant-Based Valid SA  & 93.7\%  & 100.0\%& 93.2\%  \\
\bottomrule
\multicolumn{4}{l}{\footnotesize $^\S$2 Few-Shot instances failed schema parsing; VA computed over 28 parsed responses.}
\end{tabular}
\end{table}

A superficial reading concludes that CoT verbalization causes general topological
reasoning breakdown. However, cross-examining Table~\ref{tab:wfg_taxonomy} and
Table~\ref{tab:invariant_sa} reveals a deeper, nuanced scientific truth:

\begin{enumerate}
\item \textbf{High Invariant Validity and Verdict Accuracy:} Under invariant checking,
  Zero-Shot achieves \textbf{93.7\% valid SA} (90.0\% VA), Few-Shot achieves
  \textbf{100.0\% valid SA} (100.0\% VA), and CoT achieves \textbf{93.2\% valid SA}
  (93.3\% VA). Models across all three strategies execute sound DFS traversals.

\item \textbf{Why Strict SA Collapses (The Serialization Priming Effect):} In Few-Shot
  prompting, the in-context exemplar explicitly demonstrates the solver's canonical
  tie-breaking convention (alphabetical neighbor selection). Without this formatting
  prompt, Zero-Shot and CoT models explore valid neighbors in non-canonical order. Because
  strict positional matching requires step $t$ to match identically, a single
  non-canonical branch at Step~1 invalidates every subsequent positional comparison,
  resulting in $\FEP = 1$ in 100\% of cases and depressing strict SA to 4.4\%--7.5\%.

\item \textbf{Residual CoT Tracking Drift:} Although CoT maintains 93.2\% invariant
  validity, Table~\ref{tab:wfg_taxonomy} confirms that verbalization induces subtle
  tracking drift: \texttt{visited\_set\_corruption} occurs in 100\% of CoT traces
  (vs.\ 23.3\% in Few-Shot), and false cycle hallucination rises from 0\% in Few-Shot
  to 6.7\% in CoT. In long verbalized traces, models occasionally conflate backtracked
  nodes with active recursion path elements.
\end{enumerate}

\textbf{Case Study: Localized CoT Backtrack Drift (Instance \texttt{wfg\_3ceec919261d}).}
Ground truth for this instance is \texttt{deadlock=False} (DAG). At step~9, the model's
trace correctly backtracks P4 out of the recursion stack
(\texttt{"recursion\_stack":~["P0"]}), marking P4 as fully explored. However, 12 steps
later at step~21, the model re-asserts P4 as being ``in the recursion stack'' and
declares a false cycle---contradicting its own step-9 state. This demonstrates how
unconstrained natural language tokens can induce localized state-binding drift at
terminal steps.

\subsection{First Error Position (FEP) Dynamics}

Evaluating the discrete point of first failure ($\FEP$) reveals fundamentally distinct
degradation dynamics across algorithmic domains:
\begin{itemize}
\item \textbf{Immediate Topological Divergence (WFG):} Under strict matching,
  $\FEP = 1$ in 100.0\% of error cases for ZS and CoT (mean $\FEP = 1.00$) due to
  branch tie-break divergence at the initial node. Under Few-Shot, format priming aligns
  canonical order, delaying initial failure to mean $\FEP = 8.38$ (median 7.0).
\item \textbf{Early Matrix Misalignment (Banker's):} $\FEP = 1$ occurs in 76.0\% (ZS),
  59.1\% (FS), and 82.4\% (CoT) of error instances (mean $\FEP \in [1.28, 1.55]$).
\item \textbf{Gradual Multi-Step State Drift (Buddy System):} $\FEP = 1$ occurs in only
  22.7\%--34.8\% of errors, with a mean $\FEP$ of \textbf{6.96 steps (ZS), 7.57 (FS),
  and 6.55 (CoT)} (median $\FEP = 4.0$--$5.0$). Models correctly perform initial block
  splits, but accumulate arithmetic fragmentation errors over long horizons.
\end{itemize}

\subsection{Equivalence-Aware Invariant Scoring vs.\ Strict Positional Matching}

\begin{table}[t]
\centering
\caption{Invariant-Based Valid SA vs.\ Strict Positional SA (Gemini)}
\label{tab:invariant_sa}
\begin{tabular}{L{2.4cm}C{1.5cm}C{1.5cm}C{1.4cm}}
\toprule
\textbf{Task / Variant} & \textbf{Strict SA} & \textbf{Valid SA} & \textbf{$\Delta$} \\
\midrule
\multicolumn{4}{l}{\emph{Banker's Safety ($N{=}90$)}} \\
\quad Zero-Shot    & 32.9\% & 74.4\% & $+41.6\%$ \\
\quad Few-Shot     & 41.9\% & 78.9\% & $+37.1\%$ \\
\quad CoT          & 49.7\% & 69.5\% & $+19.8\%$ \\
\quad \textbf{Mean}& \textbf{41.5\%} & \textbf{74.3\%} & $\mathbf{+32.8\%}$ \\
\midrule
\multicolumn{4}{l}{\emph{WFG Deadlock ($N{=}88$)}} \\
\quad Zero-Shot    &  4.4\% & 93.7\% & $+89.4\%$ \\
\quad Few-Shot     & 85.4\% & 100.0\%& $+14.6\%$ \\
\quad CoT          &  7.5\% & 93.2\% & $+85.8\%$ \\
\quad \textbf{Mean}& \textbf{31.2\%} & \textbf{95.6\%} & $\mathbf{+64.3\%}$ \\
\bottomrule
\end{tabular}
\end{table}

As detailed in Table~\ref{tab:invariant_sa}:
\begin{enumerate}
\item \textbf{Banker's Algorithm:} When verifying that each chosen candidate process
  satisfies $\mathbf{Need}_i \le \text{Work}$ and properly reclaims resources under the
  model's own state, valid step accuracy reaches \textbf{74.3\%} ($+32.8\%$ over
  positional SA). In 7 out of 90 instances, models executed a sound safe sequence that
  prioritized a different candidate process than the reference solver's deterministic
  tie-break.
\item \textbf{WFG Cycle Detection:} Under invariant checking (valid neighbor traversal
  and backtracking), WFG valid step accuracy reaches \textbf{95.6\%} ($+64.3\%$ over
  positional SA).
\item \textbf{Trace Length Alignment Sensitivity:} Evaluating monotonic Longest Common
  Subsequence (LCS) alignment reveals that trace length disparity
  ($m / T^* \approx 1.3$--$2.8$) creates substantial artificial cascade penalties:
  aligned SA increases Banker's CoT from 49.7\% to 63.5\% ($+13.8\%$) and WFG Few-Shot
  from 85.4\% to 97.7\% ($+12.3\%$).
\end{enumerate}

\subsection{Cross-Model Schema Fragility and Formatting Brittleness}

Table~\ref{tab:cross_model} underscores a profound divergence between proprietary
instruction-tuned models and open-weight architectures. While Gemini-3.5-Flash-Lite
achieved 99.2\% parse success, \textbf{NVIDIA Nemotron-3-Super-120B achieved only 19.7\%
parse success}, with the remaining 80.3\% of outputs failing JSON decoding due to missing
closing delimiters, extraneous conversational preambles, or unescaped tokens.

Crucially, auditing the parsed subset reveals that Nemotron achieved \textbf{95.9\%
Verdict Accuracy} and \textbf{48.5\% Step Accuracy}. Our dual-metric decomposition
($\VA_{\text{parsed}}$ vs.\ $\VA_{\text{all}}$) establishes that open-weight models
suffer not from semantic reasoning deficits, but from strict output-channel schema
adherence failures.

% ============================================================
\section{Discussion and Practical Implications}
\label{sec:discussion}
% ============================================================

\subsection{Token-Level Format Following vs.\ Algorithmic Competence}

Our findings highlight that raw parameter count does not ensure procedural fidelity.
Nemotron-120B contains 120B total parameters, yet
failed schema parsing on 80.3\% of tasks where Gemini-3.5-Flash-Lite parsed cleanly.
For practical systems verification pipelines, constrained decoding (e.g., grammar-based
sampling via outlines or JSON schemas) is essential before open-weight models can be
reliably deployed.

\subsection{Implications for Automated Verification and Systems Pedagogy}

\begin{enumerate}
\item \textbf{Never rely on unverified verdicts:} On balanced 50/50 concurrency
  distributions, LLMs exhibit a 35.8\% divergence between outcome correctness and
  procedural fidelity. Downstream systems tools must verify intermediate execution traces.
\item \textbf{Distinguish between format compliance and procedural validity:} In graph
  and state-machine problems, few-shot exemplars serve primarily to prime arbitrary
  serialization conventions (tie-breaks). Automated benchmarks must adopt equivalence-aware
  verifiers to avoid mistaking valid alternative traversals for reasoning failures.
\item \textbf{Constrain generative freedom in topological tasks:} Unstructured
  verbalization in CoT increases state-binding drift on recursion stacks. Structured
  schema emission should be enforced.
\end{enumerate}

% ============================================================
\section{Limitations}
\label{sec:limitations}
% ============================================================

This empirical investigation is bounded by several methodology considerations:
\begin{enumerate}
\item \textbf{Single-Run Deterministic Sampling:} Inferences were performed at
  $\text{temperature} = 0.0$. We did not evaluate temperature-swept self-consistency or
  majority voting across multiple runs.
\item \textbf{Per-Cell Sample Size:} With 125 instances spanning 5 tasks and 5 complexity
  levels, individual task-complexity sub-cells contain $N = 4$--$6$ instances, though
  each prompt strategy evaluates $N = 20$--$30$ instances per algorithm family.
\item \textbf{Partial Evaluation of Open Models:} Upstream API rate-limiting on
  open-router endpoints restricted open-weight evaluations (GPT-OSS-20B and LFM-2.6B)
  to WFG subsets, necessitating conditional reporting over parsed subsets.
\item \textbf{Synthetic Problem Horizons:} Instances model canonical textbook concurrency
  algorithms ($P \le 6, R \le 4$) rather than production-scale distributed systems traces.
\item \textbf{Static Prompting Regimes:} We evaluated static ZS, FS, and CoT prompts
  without dynamic tool feedback, external Python interpreter execution, or
  self-correction loops.
\item \textbf{Tie-Breaking Sensitivity:} While invariant-aware scoring resolves
  alternative valid sequences, strict positional SA remains sensitive to canonical solver
  ordering.
\end{enumerate}

% ============================================================
\section{Future Work}
\label{sec:future}
% ============================================================

Future extensions of SVAC-Concurrency will investigate:
\begin{enumerate}
\item \textbf{Grammar-Constrained Decoding:} Evaluating whether formal context-free
  grammars (CFGs) eliminate open-weight formatting collapse.
\item \textbf{Reinforcement Learning with Step Verifiers:} Utilizing SVAC's intermediate
  step verifiers as dense process reward models (PRMs) to train reasoning models on
  discrete concurrency invariants.
\item \textbf{Kernel-Level Concurrency Traces:} Extending problem instances to real Linux
  kernel eBPF trace logs and lockdep synchronization graphs.
\end{enumerate}

% ============================================================
\section{Conclusion}
\label{sec:conclusion}
% ============================================================

We presented \textbf{SVAC-Concurrency}, a fine-grained, open-source benchmark for
evaluating LLM intermediate reasoning traces on Operating System concurrency algorithms.
Through multi-dimensional metrics ($\SA$, $\FEP$, $\ECR$, and Equivalence-Aware
Validity), we showed that high outcome accuracy masks severe intermediate state failures,
deconstructed the CoT Inversion Effect in topological deadlock detection as serialization
tie-breaking priming, and demonstrated the critical impact of invariant-aware scoring.
All benchmark datasets, reference solvers, and evaluation harnesses are made publicly
available to advance verifiable algorithmic reasoning.

\balance
\bibliographystyle{IEEEtran}
\bibliography{references}

\end{document}
